% Source: https://artofproblemsolving.com/community/c3h3251973_temperature_rounding
% Stable identity: aops-topic-3251973-q01
Write $F-32=9q+r$, $0\le r\le8$. The first rounded value is $5q+\operatorname{round}(\dfrac{5r}{9})$. After returning to Fahrenheit the residual integer is $\operatorname{round}(\dfrac{9}{5}\operatorname{round}(\dfrac{5r}{9}))$. For $r=0,1,\ldots,8$, these are $0,2,2,4,4,5,5,7,7$. Thus a change occurs exactly when $F-32\equiv1,3,6,8\pmod9$, or $\boxed{F\equiv2,4,6,8\pmod9}$. Examples are $33$ and $35$ degrees Fahrenheit, which return as $34$ and $36$. Conversely, rounding an exact Fahrenheit conversion introduces an error of at most $\dfrac{2}{5}$ degrees Fahrenheit (its fractional parts are fifths). Conversion back changes the original Celsius integer by at most $\dfrac{5}{9}\cdot\dfrac{2}{5}=\dfrac{2}{9}<\dfrac{1}{2}$, so final rounding restores it. Thus no Celsius integer changes, and starting in Celsius avoids the effect. None of these calculations has a rounding tie.
